% GENERATED FILE. Edit canonical lessons, then run npm run generate:books.
% canonical-source-hash: fnv1a64:6a52b81252e00209
% canonical-lessons: AR-C285-ray, AR-C285-ataqidu-anna, AR-C285-sual, AR-W285-hamza-on-waw, AR-C285-li-dhalika-fa, AR-R285-opinion-and-so

\chapter{What I Think, And So}
\label{ch:ar-what-i-think-and-so}
\hlchaptermodality{285}

\begin{chapteropening}
\textbf{By the end of this chapter:} \emph{I can ask a question, give my opinion, and say what followed from a reason.}

Complete AR-R285-opinion-and-so: an opinion with fī raʾyī or aʿtaqidu anna, its reason with li-anna, and its result with li-dhālika or fa-.
\end{chapteropening}

\section[\texorpdfstring{\ar{رأي} (raʾy)}{رأي (raʾy)}]{\ar{رأي} (raʾy) — an opinion}

\label{lesson:AR-C285-ray}



\begin{quote}
Tell a two-step story: \textquotedblleft{}first I read, then I wrote\textquotedblright{}. (\emph{Awwalan qaraʾtu, thumma katabtu}.) Then \textquotedblleft{}I told the story when he came\textquotedblright{}. (\emph{Ḥakaytu al-qiṣṣa ʿindamā jāʾa}.)

A story says what happened. The next node asks what you \textbf{think} about it, and why.
\end{quote}

\subsection*{You'll want to know: \ar{رأي}}
\textbf{\ar{رأي}} — \emph{raʾy} — \textquotedblleft{}an opinion, a view\textquotedblright{}.

\begin{quote}
\textbf{\ar{ما} \ar{رأيك} \ar{في} \ar{الفيلم؟}} — \emph{mā raʾyuka fī al-fīlm?} — What do you think of the film?
\end{quote}

\begin{quote}
\textbf{\ar{في} \ar{رأيي،} \ar{الفيلم} \ar{مفيد}.} — \emph{fī raʾyī, al-fīlm mufīd} — In my opinion, the film is useful.
\end{quote}

\emph{Raʾyī} is \textquotedblleft{}my opinion\textquotedblright{}, with the \emph{-ī} of \textquotedblleft{}my\textquotedblright{}. \emph{Raʾyuka} is \textquotedblleft{}your opinion\textquotedblright{} to a man; to a woman it is \emph{raʾyuki}. \emph{Fī}, \textquotedblleft{}in\textquotedblright{}, says what the opinion is about.

\begin{cousinweb}
\textbf{\ar{رأي}} (\emph{raʾy}) is built on the native Arabic root \emph{r-a-y}, \textquotedblleft{}to see\textquotedblright{} (Lane): the root of \emph{raʾā}, \textquotedblleft{}he saw\textquotedblright{}, from chapter twenty-nine. Lane gives \emph{raʾā fī al-amr raʾyan} for forming a view of a matter. An opinion, in Arabic, is what you see in something.
\end{cousinweb}

\subsection*{Your turn}
Say these aloud:

\begin{itemize}
\raggedright
  \item \emph{raʾy}, with the catch in the middle — \emph{raʾ-y}
  \item \textquotedblleft{}What do you think of the film?\textquotedblright{} — \emph{mā raʾyuka fī al-fīlm?}
  \item \textquotedblleft{}In my opinion, the book is useful\textquotedblright{} — \emph{fī raʾyī, al-kitāb mufīd}
  \item the same question to a woman, about a sentence — \emph{mā raʾyuki fī al-jumla?}
\end{itemize}

\begin{tcolorbox}[breakable,colback=teal!4,colframe=teal!35!black,title={Before you move on}]
What does \ar{رأي} mean? (\textquotedblleft{}An opinion\textquotedblright{}.) \textquotedblleft{}In my opinion\textquotedblright{}? (\emph{\textbf{Fī raʾyī.}}) Which verb shares its root? (\emph{\textbf{Raʾā}}, he saw.)
\end{tcolorbox}
\section[\texorpdfstring{\ar{أعتقد} \ar{أن} (aʿtaqidu anna)}{أعتقد أن (aʿtaqidu anna)}]{\ar{أعتقد} \ar{أن} (aʿtaqidu anna) — I think that}

\label{lesson:AR-C285-ataqidu-anna}



\begin{quote}
Say \textquotedblleft{}in my opinion\textquotedblright{}. (\emph{Fī raʾyī}.) Then say \emph{aktub anna Sāmī jārī}, \textquotedblleft{}I write that Sami is my neighbour\textquotedblright{}.

\emph{Anna} hangs a whole sentence on a verb. Hang it on a verb of thinking, and you can give an opinion.
\end{quote}

\begin{grammarlens}[title={Grammar Lens: aʿtaqidu anna, aẓunnu anna}]
\emph{Iʿtaqada}, to believe, from chapter one hundred and sixty-four, has the present \emph{aʿtaqidu}, I believe, I think. Put \emph{anna} after it:

\begin{quote}
\textbf{\ar{أعتقد} \ar{أن} \ar{الامتحان} \ar{صعب}.} — \emph{aʿtaqidu anna al-imtiḥāna ṣaʿb} — I think the exam is hard.
\end{quote}

The noun after \emph{anna} takes \emph{-a}: \emph{al-imtiḥāna}. With \textquotedblleft{}it\textquotedblright{} the pronoun joins \emph{anna} itself: \emph{aʿtaqidu annahu ṣaʿb}, I think it's hard.

\textbf{\ar{أظن} \ar{أن}} — \emph{aẓunnu anna} — is \textquotedblleft{}I suppose\textquotedblright{}, a little less sure: \emph{aẓunnu anna al-qiṭāra sarīʿ}, I suppose the train is fast.
\end{grammarlens}

\begin{grammarlens}[title={Grammar Lens: agreeing and not}]
For \textquotedblleft{}I don't think so\textquotedblright{}, put \emph{lā} in front and \emph{dhālika}, \textquotedblleft{}that\textquotedblright{}, after: \textbf{\ar{لا} \ar{أعتقد} \ar{ذلك}} — \emph{lā aʿtaqidu dhālika}.

\begin{quote}
— \emph{mā raʾyuka fī al-fīlm?} — \emph{aʿtaqidu annahu mumtāz.} — \emph{lā aʿtaqidu dhālika!}
\end{quote}

Three ways to start an opinion, then: \emph{fī raʾyī}, \emph{aʿtaqidu anna}, \emph{aẓunnu anna}. English drops its \textquotedblleft{}that\textquotedblright{}; Arabic keeps \emph{anna}.
\end{grammarlens}

\subsection*{Your turn}
Say these aloud:

\begin{itemize}
\raggedright
  \item \textquotedblleft{}I think the exam is hard\textquotedblright{} — \emph{aʿtaqidu anna al-imtiḥāna ṣaʿb}
  \item \textquotedblleft{}I think it's excellent\textquotedblright{} — \emph{aʿtaqidu annahu mumtāz}
  \item \textquotedblleft{}I suppose the train is fast\textquotedblright{} — \emph{aẓunnu anna al-qiṭāra sarīʿ}
  \item \textquotedblleft{}I don't think so\textquotedblright{} — \emph{lā aʿtaqidu dhālika}
  \item \textquotedblleft{}I think the lawyer is here\textquotedblright{} — \emph{aʿtaqidu anna al-muḥāmiya hunā}
\end{itemize}

\begin{tcolorbox}[breakable,colback=teal!4,colframe=teal!35!black,title={Before you move on}]
\textquotedblleft{}I think that\textquotedblright{}? (\emph{\textbf{Aʿtaqidu anna.}}) The less sure one? (\emph{\textbf{Aẓunnu anna.}}) What happens to the noun after \emph{anna}? (\textbf{It takes \emph{-a}}: \emph{al-imtiḥāna}.) \textquotedblleft{}I don't think so\textquotedblright{}? (\emph{\textbf{Lā aʿtaqidu dhālika.}})
\end{tcolorbox}
\section[\texorpdfstring{\ar{سؤال} (suʾāl)}{سؤال (suʾāl)}]{\ar{سؤال} (suʾāl) — a question}

\label{lesson:AR-C285-sual}



\begin{quote}
Say \textquotedblleft{}I think that\textquotedblright{}, then \textquotedblleft{}he asked\textquotedblright{}. (\emph{Aʿtaqidu anna}, \emph{saʾala}.)

Before an opinion there is often a question. Here is the noun for it.
\end{quote}

\subsection*{You'll want to know: \ar{سؤال}}
\textbf{\ar{سؤال}} — \emph{suʾāl} — \textquotedblleft{}a question\textquotedblright{}.

\begin{quote}
\textbf{\ar{عندي} \ar{سؤال}.} — \emph{ʿindī suʾāl} — I have a question.
\end{quote}

\emph{ʿIndī}, \textquotedblleft{}at me\textquotedblright{}, is how Arabic says \textquotedblleft{}I have\textquotedblright{}. The answer to a question is \emph{jawāb}, from chapter seventy-four, and the two make a pair: \emph{suʾāl wa-jawāb}, question and answer.

Say it \emph{su-ʾāl}, with the catch before the long \emph{ā}: the same catch as in the middle of \emph{saʾala}.

\begin{cousinweb}
\textbf{\ar{سؤال}} (\emph{suʾāl}) is built on the native Arabic root \emph{s-a-l}, \textquotedblleft{}to ask\textquotedblright{} (Lane), the root of \emph{saʾala}, \textquotedblleft{}he asked\textquotedblright{}. Lane gives \emph{suʾāl} as \textquotedblleft{}a question\textquotedblright{}, the counterpart of \emph{jawāb}. The catch in the middle is the root's own middle letter, a hamza. In \emph{saʾala} it sits on \textbf{\ar{ا}}. Here the vowel before it is \emph{u}, and it is written on a different seat, which the next lesson puts in your hand.
\end{cousinweb}

\subsection*{Your turn}
Say these aloud:

\begin{itemize}
\raggedright
  \item \emph{suʾāl}, slowly, with the catch — \emph{su-ʾāl}
  \item \textquotedblleft{}I have a question\textquotedblright{} — \emph{ʿindī suʾāl}
  \item \textquotedblleft{}question and answer\textquotedblright{} — \emph{suʾāl wa-jawāb}
  \item \textquotedblleft{}I think the question is easy\textquotedblright{} — \emph{aʿtaqidu anna as-suʾāla sahl}
\end{itemize}

\begin{tcolorbox}[breakable,colback=teal!4,colframe=teal!35!black,title={Before you move on}]
What does \ar{سؤال} mean? (\textquotedblleft{}A question\textquotedblright{}.) Which verb shares its root? (\emph{\textbf{Saʾala}}, he asked.) What is its partner? (\emph{\textbf{Jawāb}}, an answer.)
\end{tcolorbox}
\section[\texorpdfstring{\ar{ؤ} (hamza on waw)}{ؤ (hamza on waw)}]{\ar{ؤ} (hamza on waw) — the catch inside \ar{سؤال}}

\label{lesson:AR-W285-hamza-on-waw}



\begin{quote}
Before the new one: write \ar{ء} and \ar{و} once each, and name them.

Now say \textbf{\ar{سؤال}} (\emph{suʾāl}), \textquotedblleft{}a question\textquotedblright{}. You met it a moment ago. This lesson takes its catch out and puts it in your hand.
\end{quote}

\subsection*{Script you'll notice: \ar{ؤ}}
\textbf{\ar{ؤ}} — \emph{hamza} on a \emph{wāw} seat, the glottal catch written \emph{ʾ} in this book's romanization.

It is the second letter of \textbf{\ar{سؤال}}, between \textbf{\ar{س}} and \textbf{\ar{ا}}. It is two pieces: a \textbf{\ar{و}} below, the seat, and the small hamza mark riding on top of it. Like a plain \textbf{\ar{و}}, it takes the hand of \textbf{\ar{س}} on its right and does not join to the letter after it, so a gap opens before \textbf{\ar{ال}}.

The seat does not change the sound. \textbf{\ar{أ}}, \textbf{\ar{ء}} and \textbf{\ar{ؤ}} all say the same catch; the seat follows the vowels around it, and here the \emph{u} before it calls for \textbf{\ar{و}}.

\subsection*{Writing: \ar{ؤ}}
Put your pen on \ar{ؤ} and follow its line. Write the \textbf{\ar{و}} first, then add the hamza above it. Copy the shape you can see — slowly, and larger than it is printed.

\begin{quote}
This book does not yet tell you \textbf{where to start this character or which way to travel}. It is not written down here until it can be written down with a source. Copying what is in front of you needs no such source, and it is how the shape gets into your hand in the meantime.
\end{quote}

\subsection*{Your turn}
\begin{itemize}
\raggedright
  \item \emph{Find:} the seat and the mark of \ar{ؤ} inside \textbf{\ar{سؤال}}
  \item \emph{Trace it:} \ar{ؤ} three times, saying the catch as you finish each one
  \item \emph{Write it:} \ar{سؤ}, joined, then the gap, then \ar{ال}
  \item \emph{Write it:} \textbf{\ar{سؤال}} whole, then \textbf{\ar{سأل}} beside it: one catch, two seats
  \item \emph{Read it:} \textbf{\ar{سؤال}} aloud once more
\end{itemize}

\begin{tcolorbox}[breakable,colback=teal!4,colframe=teal!35!black,title={Before you move on}]
What is this — \ar{ؤ}? (\emph{\textbf{Hamza}} on a \emph{wāw} seat, a catch.) Which word did it come from? (\textbf{\ar{سؤال}}, \emph{suʾāl}, \textquotedblleft{}a question\textquotedblright{}.) Does the seat change the sound? (\textbf{No}: only the spelling.)
\end{tcolorbox}
\section[\texorpdfstring{\ar{لذلك،} \ar{فـ} (li-dhālika, fa-)}{لذلك، فـ (li-dhālika, fa-)}]{\ar{لذلك،} \ar{فـ} (li-dhālika, fa-) — so}

\label{lesson:AR-C285-li-dhalika-fa}



\begin{quote}
Say \textquotedblleft{}I have a question\textquotedblright{}. (\emph{ʿIndī suʾāl}.) Then \textquotedblleft{}because\textquotedblright{}. (\emph{Li-anna}.)

\emph{Li-anna} puts the reason after the claim. A story often wants the opposite order: first the reason, then what came of it.
\end{quote}

\begin{grammarlens}[title={Grammar Lens: li-dhālika, after the reason}]
\begin{quote}
\textbf{\ar{نمت} \ar{مبكرا} \ar{لأنني} \ar{كنت} \ar{متعبا}.} — \emph{nimtu mubakkiran li-annanī kuntu mutʿaban} — I went to bed early because I was tired.
\end{quote}

\begin{quote}
\textbf{\ar{كنت} \ar{متعبا،} \ar{لذلك} \ar{نمت} \ar{مبكرا}.} — \emph{kuntu mutʿaban, li-dhālika nimtu mubakkiran} — I was tired, so I went to bed early.
\end{quote}

The same two facts. With \emph{li-anna} the claim comes first; with \emph{li-dhālika}, \textquotedblleft{}therefore, so\textquotedblright{}, the reason comes first and the result follows. The \emph{-an} on \emph{mutʿaban} after \emph{kuntu} is \emph{kāna}'s doing; chapter two hundred and eighty-nine comes back to it.
\end{grammarlens}

\begin{grammarlens}[title={Grammar Lens: fa-, the short \textquotedblleft{}so\textquotedblright{}}]
\textbf{\ar{فـ}} — \emph{fa-} — does the same job in one letter. Like \emph{wa-}, it is written onto the next word, with no space:

\begin{quote}
\textbf{\ar{تأخر} \ar{القطار،} \ar{فأخذنا} \ar{الحافلة}.} — \emph{taʾakhkhara al-qiṭār, fa-akhadhnā al-ḥāfila} — The train was late, so we took the bus.
\end{quote}

\emph{Wa-} adds the next thing; \emph{fa-} says the next thing came of the first. \emph{Li-dhālika} is the heavier of the two, and it suits a reason that needs weight: \emph{kāna al-imtiḥān ṣaʿban, li-dhālika darastu kathīran}, the exam was hard, so I studied a lot.
\end{grammarlens}

\subsection*{Your turn}
Say these aloud:

\begin{itemize}
\raggedright
  \item \emph{nimtu mubakkiran li-annanī kuntu mutʿaban}
  \item the same facts the other way — \emph{kuntu mutʿaban, li-dhālika nimtu mubakkiran}
  \item \textquotedblleft{}The train was late, so we took the bus\textquotedblright{} — \emph{taʾakhkhara al-qiṭār, fa-akhadhnā al-ḥāfila}
  \item an opinion and its result — \emph{aʿtaqidu anna al-imtiḥāna ṣaʿb, li-dhālika adrusu}
\end{itemize}

\begin{tcolorbox}[breakable,colback=teal!4,colframe=teal!35!black,title={Before you move on}]
After \emph{li-anna}, what comes? (\textbf{The reason.}) After \emph{li-dhālika}? (\textbf{The result.}) The one-letter \textquotedblleft{}so\textquotedblright{}? (\textbf{\ar{فـ}}, \emph{fa-}, written onto the next word.) What do \emph{li-anna} and \emph{li-dhālika} share? (\emph{\textbf{Li-}}, for.)
\end{tcolorbox}
\section[(review)]{Review — an opinion, its reason and what came of it}

\label{lesson:AR-R285-opinion-and-so}



\begin{quote}
Retell the trip to the sea in three steps, with \emph{awwalan}, \emph{thumma} and \emph{fī an-nihāya}. Then say how it ended, with \emph{ʿindamā}: \emph{ʿindamā waṣaltu, ḥakaytu al-qiṣṣa}. And one step with \emph{baʿda an}: \emph{baʿda an akalnā, sabaḥnā}.
\end{quote}

\begin{grammarlens}[title={What you've built}]
Four lessons, and an opinion has its parts. A question to open it: \textbf{\ar{عندي} \ar{سؤال}}, \emph{ʿindī suʾāl}. Three ways to start it: \emph{fī raʾyī}, \emph{aʿtaqidu anna}, \emph{aẓunnu anna}. A reason after the claim, with \emph{li-anna}; a result after the reason, with \emph{li-dhālika} or \emph{fa-}.

All of it in one breath:

\begin{quote}
\textbf{\ar{في} \ar{رأيي} \ar{القصة} \ar{جميلة،} \ar{لذلك} \ar{أحكيها} \ar{كثيرا}.} — \emph{fī raʾyī al-qiṣṣa jamīla, li-dhālika aḥkīhā kathīran} — In my opinion the story is lovely, so I tell it often.
\end{quote}
\end{grammarlens}

\subsection*{Your turn}
Say these aloud:

\begin{itemize}
\raggedright
  \item \textquotedblleft{}What do you think of the story?\textquotedblright{} — \emph{mā raʾyuka fī al-qiṣṣa?}
  \item \textquotedblleft{}I think the number is right\textquotedblright{} — \emph{aʿtaqidu anna ar-raqma ṣaḥīḥ}
  \item \textquotedblleft{}The information was useful, so I wrote it\textquotedblright{} — \emph{kānat al-maʿlūmāt mufīdatan, fa-katabtuhā}
  \item an opinion of your own, its reason with \emph{li-anna}, and a result with \emph{li-dhālika}
\end{itemize}

\begin{tcolorbox}[breakable,colback=teal!4,colframe=teal!35!black,title={Before you move on}]
\textquotedblleft{}In my opinion\textquotedblright{}? (\emph{\textbf{Fī raʾyī.}}) \textquotedblleft{}I have a question\textquotedblright{}? (\emph{\textbf{ʿIndī suʾāl.}}) Which seat carries the hamza in \textbf{\ar{سؤال}}? (\textbf{\ar{و}.}) The one-letter \textquotedblleft{}so\textquotedblright{}? (\emph{\textbf{Fa-.}})
\end{tcolorbox}
